\thispagestyle{plain}

\section*{Abstract}
\label{sec:abstract}

Invasive alien plants are among the main drivers of biodiversity loss,
and controlling them requires up-to-date maps of where they grow.
Satellite imagery can provide such maps, but the supervised classifiers
it relies on need reference data, locations where the species is known
to be present, and collecting these in the field is slow and costly.
Existing distribution products are updated infrequently, are spatially
coarse and rarely reach the species level.

Citizen-science platforms such as iNaturalist hold thousands of free,
georeferenced observations of invasive plants. Rich as they are, these
data are isolated points of uneven positional and taxonomic quality,
concentrated where observers happen to go, and they record presences
only, never confirmed absences. Using them directly to train a
satellite-based classifier is therefore not straightforward.

Earth Observation foundation models, pre-trained on large volumes of
satellite data, offer a possible way through. They turn each pixel into
a general-purpose representation, an embedding, on which a light
classifier can be trained from limited and noisy labels, potentially
making citizen-science observations usable as reference data.

This intermediate report, developed within the MAP.INVADER project,
investigates to what extent citizen-science observations, turned into
reference data and represented by foundation-model embeddings, can
reduce the dependency on costly expert reference data for mapping
\emph{Acacia dealbata} in the Coimbra district, Portugal. It reviews
the use of citizen science and foundation models in Earth Observation,
proposes a pipeline that converts point observations into training
labels for satellite imagery, and sets out an evaluation that measures
how much is gained by adding expert data to a citizen-science base.

\section*{Keywords}
\label{sec:keywords}

Invasive alien plants; citizen science; Earth Observation foundation
models; satellite imagery; species mapping; reference data.
